\section{A priori bounds}\label{sec:apriori}

Throughout this section $f$ satisfies $0<\alpha\le f(k)\le\beta<\infty$ for all $k\ge0$, and $C=\beta/\alpha$.
The $m$ edges sent by a vertex $t$ are called the \emph{slots} of $t$ and are numbered $1,\dots,m$ in the order
in which their endpoints are chosen.

\begin{lemma}[Joint edge bound]\label{lem:edges}
Let $E$ be a finite multiset of pairs $\{i,j\}$ of distinct labels in $[n]$, with $e$ elements counted with
multiplicity. The probability that $G_n$ contains every pair of $E$ as an edge, with at least the prescribed
multiplicity, is at most
\[
   \prod_{\{i,j\}\in E}\frac{mC}{\max(i,j)-1}\;\le\;(2mC)^e\prod_{\{i,j\}\in E}\frac1{\max(i,j)} .
\]
\end{lemma}

\begin{proof}
At any slot of a vertex $t$ and for any history, a given vertex $i<t$ is chosen with conditional probability
$f(D_i)/S\le\beta/(\alpha(t-1))=C/(t-1)$, because the denominator is a sum of $t-1$ weights, each at least $\alpha$.
Assign to every element $\{i,j\}$ of $E$ a slot of its younger endpoint $\max(i,j)$, distinct elements receiving
distinct slots. Exposing the slots in chronological order and conditioning successively on the past, the
probability that every assigned slot chooses its prescribed endpoint is at most the product of $C/(\max(i,j)-1)$
over $E$. There are at most $m^e$ assignments, and the event in the statement is the union of the corresponding
events. The second inequality uses $t-1\ge t/2$ for $t\ge2$.
\end{proof}

\begin{lemma}\label{lem:schur}
The symmetric $n\times n$ matrix $H=(1/\max(i,j))_{i,j\le n}$ has operator norm at most $4$ on $\ell^2$.
\end{lemma}

\begin{proof}
Let $w_i=i^{-1/2}$. Then $\sum_{j\le i}H_{ij}w_j=i^{-1}\sum_{j\le i}j^{-1/2}\le 2i^{-1/2}$ and
$\sum_{j>i}H_{ij}w_j=\sum_{j>i}j^{-3/2}\le 2i^{-1/2}$, so $Hw\le 4w$ coordinatewise. For real $x$, the inequality
$2|x_ix_j|\le x_i^2w_j/w_i+x_j^2w_i/w_j$ and the symmetry of $H$ give
$|x^{\mathsf T}Hx|\le\sum_i x_i^2 (Hw)_i/w_i\le4\|x\|^2$.
\end{proof}

Let $U_n$ be uniform on $[n]$ and independent of $G_n$, and let $B_r(G_n,U_n)$ be the ball of radius $r$ around
$U_n$ for the graph distance that ignores orientation. Put $B=8mC$.

\begin{proposition}[Balls are bounded in mean and avoid old vertices]\label{prop:balls}
For all $n\ge1$, $r\ge0$ and $1\le K\le n$,
\[
  \E\,|B_r(G_n,U_n)|\le\sum_{l=0}^{r}B^l,\qquad
  \Prob\bigl(B_r(G_n,U_n)\cap[K]\neq\emptyset\bigr)\le\Bigl(\sum_{l=0}^{r}B^l\Bigr)\sqrt{K/n}.
\]
\end{proposition}

\begin{proof}
A vertex at distance $l$ from $U_n$ is the endpoint of a path $v_0=U_n,v_1,\dots,v_l$ with distinct vertices,
hence distinct edges. By Lemma~\ref{lem:edges}, the expected number of such paths ending in a set $A$ is at most
\[
  \frac{(2mC)^l}{n}\sum_{v_0,\dots,v_l}\prod_{k=1}^{l}\frac{1}{\max(v_{k-1},v_k)}\mathbf 1_{v_l\in A}
  =\frac{(2mC)^l}{n}\,\mathbf 1^{\mathsf T}H^l\mathbf 1_A\le\frac{(2mC)^l}{n}\,\sqrt n\,4^l\sqrt{|A|},
\]
by Lemma~\ref{lem:schur}. Take $A=[n]$ and $A=[K]$ and sum over $l\le r$.
\end{proof}

\begin{proposition}[No short cycles]\label{prop:cycles}
For every fixed $r$, the probability that $B_r(G_n,U_n)$ contains a cycle or a multiple edge tends to $0$.
\end{proposition}

\begin{proof}
Fix $\eps\in(0,1)$ and let $K=\lfloor\eps n\rfloor$. Suppose that the ball avoids $[K]$ and is not a simple tree.
Take a breadth-first spanning tree of the ball and an edge of the ball outside it (a second copy of a multiple edge
counts as such an edge). This edge together with the tree paths from its endpoints to the root forms a connected
rooted multigraph with $v\le 2r+1$ vertices and exactly $v$ edges, all of whose labels exceed $K$. There are
finitely many such rooted shapes. For a fixed shape there are at most $n^v$ labellings; the root label equals
$U_n$ with probability $1/n$; and by Lemma~\ref{lem:edges} the $v$ edges are present with probability at most
$(2mC/(\eps n))^v$. The union bound gives $O_{r,\eps}(1/n)$. With Proposition~\ref{prop:balls},
\[
  \limsup_{n\to\infty}\Prob\bigl(B_r(G_n,U_n)\text{ is not a simple tree}\bigr)\le\Bigl(\sum_{l\le r}B^l\Bigr)\sqrt\eps ,
\]
and $\eps$ is arbitrary.
\end{proof}

\begin{corollary}\label{cor:tight}
The laws of the rooted graphs $(G_n,U_n)$ are tight for the local topology. Every subsequential limit is almost
surely a locally finite simple tree in which every vertex has out-degree $m$; in particular it is infinite.
\end{corollary}

\begin{proof}
A vertex has at most $m$ outgoing edges, so there are finitely many rooted directed multigraphs with at most $M$
vertices that can occur as balls. Tightness follows from the first bound of Proposition~\ref{prop:balls} and
Markov's inequality, applied for each $r$ with a threshold $M_r$ chosen so that the exceptional probabilities are
summable. Limits are simple trees by Proposition~\ref{prop:cycles}. Every vertex other than vertex $1$ has
out-degree $m$, all edges leaving $B_r$ are seen in $B_{r+1}$, and $B_{r+1}$ contains vertex $1$ with probability
at most $(\sum_{l\le r+1}B^l)n^{-1/2}$; hence all out-degrees in a limit equal $m$. A finite tree with $v$
vertices has $v-1<mv$ edges, so the limit is infinite.
\end{proof}

\begin{lemma}[Degrees]\label{lem:degrees}
Let $D_i(n)$ be the in-degree of $i$ in $G_n$. Then $\E D_{U_n}(n)^2\le mC+2(mC)^2$, and for every $\delta>0$,
\[
   \Prob\Bigl(\max_{i\le n}D_i(n)\ge\delta n\Bigr)\le\frac{mC+2(mC)^2}{\delta^2n}.
\]
\end{lemma}

\begin{proof}
Let $L_i=mC\sum_{s=i}^{n-1}1/s$. Summing the single-slot bound $C/(t-1)$ over the slots of $t=i+1,\dots,n$ gives
$\E D_i(n)\le L_i$, and summing the product of two such bounds over ordered pairs of distinct slots gives
$\E[D_i(n)(D_i(n)-1)]\le L_i^2$. Exchanging the order of summation,
\[
  \sum_{i\le n}L_i=mC(n-1),\qquad \sum_{i\le n}L_i^2=(mC)^2\sum_{s,t<n}\frac{\min(s,t)}{st}\le 2(mC)^2(n-1).
\]
Hence $\sum_i\E D_i(n)^2\le n\,(mC+2(mC)^2)$, which gives the first claim, and the second follows from Markov's
inequality and a union bound.
\end{proof}

The second statement says that no vertex receives a positive fraction of the edges. The first gives uniform
integrability of the degree of the root, so the root of every subsequential limit has mean in-degree $m$.
